%% Plan de séquence (document enseignant) : vue d'ensemble de la séquence 4,
%% séances, documents utilisés et évaluations (papier, QCM AMC, salle informatique).
%% Découpage en séances proposé : à ajuster au rythme de la classe.

%% Font size %%
\documentclass[11pt]{article}

%% Load the custom package
\usepackage{Mathdoc}

%% Numéro de séquence %% Titre de la séquence %%
\renewcommand{\centerhead}{S04 - Plan de séquence}

%% Spacing commands %%
\renewcommand{\baselinestretch}{1} \setlength{\parindent}{0pt}

\newcommand{\champ}[2]{\textbf{#1} & #2 \\ \hline}
\newcommand{\doc}[1]{\texttt{\small #1}}

\begin{document}

\begin{center}
{\Large\bfseries Séquence 4 : Les angles}\\[1mm]
{\large 6\textsuperscript{e} -- Espace et géométrie}
\end{center}

\renewcommand{\arraystretch}{1.35}
\begin{tabularx}{\linewidth}{|>{\raggedright\arraybackslash}p{3.4cm}|X|}
\hline
\champ{Objectif}{Utiliser le rapporteur avec précision (progression 6\textsuperscript{e}).}
\champ{Programme (cycle 3)}{Reconnaître et nommer un angle droit, aigu, obtus ; \textbf{comparer des angles sans avoir recours à leur mesure} (superposition, gabarit) ; mesurer et construire un angle de mesure donnée en degrés.}
\champ{Compétences}{\textbf{GEO10} Nommer un angle -- \textbf{GEO11} Sommet et côtés -- \textbf{GEO12} Nature d'un angle -- \textbf{GEO13} Classer des angles -- \textbf{GEO14} Mesurer au rapporteur -- \textbf{GEO15/16} Construire un angle de mesure donnée. \newline
\textit{GEO17 (bissectrice) : non traitée dans cette séquence.}}
\champ{Prérequis}{Point, demi-droite et sa notation $[AB)$, points alignés (séquence 2) ; angle droit et son codage.}
\champ{Fil rouge}{Calcul mental en début de séance (compléments à $100$ et à $360$, diaporama \doc{Act-6C04-03}).}
\end{tabularx}

\section*{Déroulement prévu}

\renewcommand{\arraystretch}{1.3}
\begin{tabularx}{\linewidth}{|c|>{\raggedright\arraybackslash}X|>{\centering\arraybackslash}p{1.9cm}|>{\raggedright\arraybackslash}p{5.6cm}|}
\hline
\textbf{Séance} & \textbf{Contenu} & \textbf{Comp.} & \textbf{Documents} \\ \hline
1 & \textbf{Découverte} : l'angle comme ouverture (portes, champs de vision), puis définition, sommet et côtés, nommer un angle. & GEO10, 11 &
\doc{Act-6C04-01} (découverte) ; Cours I ; \doc{Ex-6C04-01}, \doc{Ex-6C04-02} \\ \hline
2 & \textbf{Manipulation} : comparer des angles avec des gabarits (pliage d'une A5), puis angles particuliers (nul, aigu, droit, obtus, plat). & GEO12, 13 &
\doc{Act-6C04-02} (gabarits) ; Cours II ; \doc{Ex-6C04-03}, \doc{Ex-6C04-04} ; préparation \doc{Prep-6C04-01} \\ \hline
3 & \textbf{Évaluation écrite 1} (35 min), puis reprise : classer des angles. & GEO10, 11 &
\doc{DS-6C04-01} ; \doc{Ex-6C04-04} \\ \hline
4 & Mesurer un angle au rapporteur (graduation intérieure / extérieure, vérification par la nature). & GEO14 &
Cours III ; \doc{Ex-6C04-05} (ex. 1) \\ \hline
5 & Construire un angle de mesure donnée. & GEO15, 16 &
Cours IV ; \doc{Ex-6C04-05} (ex. 2), \doc{Ex-6C04-06} \\ \hline
6 & \textbf{Évaluation écrite 2} (35 min), puis entraînement : mesurer et construire. & GEO12, 13 &
\doc{DS-6C04-02} \\ \hline
7 & \textbf{Salle informatique} : bilan sur Capytale (nommer, mesurer, construire\ldots). & GEO10 à 16 &
Activité Capytale (accès par l'ENT) \\ \hline
8 & \textbf{Salle informatique} : permis rapporteur, noté sur 20. & GEO14, 15/16 &
Permis rapporteur (accès par l'ENT) \\ \hline
9 & \textbf{Évaluation écrite 3} (35 min) : mesurer et construire. & GEO14, 15/16 &
\doc{DS-6C04-03} \\ \hline
10 & \textbf{Bilan de la séquence} : QCM corrigé par AMC. & GEO10 à 13 &
\doc{DS-6C04-05} \\ \hline
\end{tabularx}

\section*{Évaluations}

\renewcommand{\arraystretch}{1.35}
\begin{tabularx}{\linewidth}{|>{\raggedright\arraybackslash}p{3.3cm}|>{\raggedright\arraybackslash}X|>{\raggedright\arraybackslash}p{3.6cm}|>{\raggedright\arraybackslash}p{3.2cm}|}
\hline
\textbf{Évaluation} & \textbf{Compétences évaluées} & \textbf{Format} & \textbf{Notation} \\ \hline
\doc{DS-6C04-01} & GEO10 nommer un angle (recto), GEO11 sommet et côtés (verso) & Papier, recto-verso, 35 min & sur 40 + grille de compétences \\ \hline
\doc{DS-6C04-02} & GEO12 nature d'un angle (recto), GEO13 classer des angles (verso) & Papier, recto-verso, 35 min & sur 40 + grille de compétences \\ \hline
\doc{DS-6C04-03} & GEO14 mesurer (recto), GEO15/16 construire (verso) & Papier, recto-verso, 35 min, rapporteur & sur 40 + grille de compétences \\ \hline
Bilan Capytale & Nommer, mesurer, construire\ldots{} (GEO10 à GEO16) & Sur ordinateur (ENT) & bilan formatif \\ \hline
Permis rapporteur & Mesurer et construire un angle (GEO14, GEO15/16) & Sur ordinateur (ENT) & sur 20 \\ \hline
\doc{DS-6C04-05} QCM & GEO10 à GEO13 : nommer, sommet et côtés, nature, classer & QCM papier, corrigé par AMC & automatique \\ \hline
\end{tabularx}


\end{document}

%%% Local Variables:
%%% mode: LaTeX
%%% TeX-master: t
%%% End:
